% Prostudováno: ucebnice/02_Funkce_a_rovnice/Zaklady_funkci/Funkce.tex (D/H, graf a vlastnosti),
% ucebnice/02_Funkce_a_rovnice/Linearni_funkce_a_rovnice/Lineární funkce.tex (průsečíky a graf),
% ucebnice/02_Funkce_a_rovnice/Linearni_funkce_a_rovnice/Linerání rovnice.tex, úvod k rovnici.
% Realisticky.cz: externi_zdroje/realisticky/02_Funkce_a_rovnice/02_Linearni_rovnice_a_nerovnice/020209_Graficke_reseni_rovnic_a_nerovnic/020209_Lekce.pdf;
% využita vazba rovnosti dvou výrazů a průsečíku dvou grafů. Úlohy a graf vlastní.
\documentclass[11pt,a4paper]{article}
% Spolecny styl pro ucitelske pripravy MAT (Work/MAT)
% Verze 1.1, 2026-10-23
\NeedsTeXFormat{LaTeX2e}
\ProvidesPackage{priprava}[2026/10/23 v1.1 Sablona ucitelskych priprav MAT]

% Zakladni typografie a cestina
\RequirePackage[T1]{fontenc}
\RequirePackage[utf8]{inputenc}
\RequirePackage[czech]{babel}
\RequirePackage{lmodern}
\RequirePackage{microtype}
\RequirePackage{amsmath,amssymb}
\RequirePackage{xcolor}
\RequirePackage{enumitem}
\RequirePackage[most]{tcolorbox}
\RequirePackage{tabularx,booktabs,array}
\RequirePackage{fancyhdr}
\RequirePackage{needspace}
\RequirePackage{geometry}
\RequirePackage[hidelinks]{hyperref}

\geometry{a4paper,top=17mm,bottom=18mm,left=23mm,right=23mm}
\definecolor{notegray}{gray}{0.39}
\setlength{\parindent}{0pt}
\setlength{\parskip}{2.5pt plus 1pt minus 0.5pt}
\setlist[itemize]{topsep=2pt,itemsep=1pt,parsep=0pt,leftmargin=1.45em}
\setlist[enumerate]{topsep=2pt,itemsep=2pt,parsep=0pt,leftmargin=1.65em}
\setcounter{secnumdepth}{0}
\renewcommand{\arraystretch}{1.12}

% Zahlavi: vlevo rocnik | globalni tema, vpravo cislo: nazev hodiny
% Zapati zustava zcela prazdne.
\pagestyle{fancy}
\fancyhf{}
\renewcommand{\headrulewidth}{0.35pt}
\renewcommand{\footrulewidth}{0pt}
\setlength{\headheight}{14pt}
\newcommand{\nastavPripravu}[4]{%
  \fancyhead[L]{\small #1\enspace\textbar\enspace #2}%
  \fancyhead[R]{\small #3: #4}%
}

% Minimalisticke cernobile boxy. Kazda definice zvlast.
\tcbset{commonlesson/.style={enhanced,breakable=false,colback=white,colframe=black,
  colbacktitle=white,coltitle=black,boxrule=0.45pt,
  arc=0.8mm,outer arc=0.8mm,left=3mm,right=3mm,top=1.6mm,bottom=1.6mm,
  toptitle=1mm,bottomtitle=0.7mm,fonttitle=\bfseries,
  attach boxed title to top left={xshift=2.5mm,yshift=-1.15mm},
  boxed title style={colback=white,colframe=white,boxrule=0pt,left=1mm,right=1mm}}}
\newtcolorbox{zadaniBox}[1][]{commonlesson,title={Zadání},#1}
\newtcolorbox{prikladBox}[1][]{commonlesson,title={Příklad},#1}
\newtcolorbox{definiceBox}[1][]{commonlesson,title={Definice},#1}

% Spolecny box pro zadani a jeho reseni, bez mezery mezi castmi.
% Pouziti:
%   \begin{prikladReseniBox}
%     Zadani ulohy.
%     \tcblower
%     \begin{align*} ... \end{align*}
%   \end{prikladReseniBox}
% Spodni cast ma jemny sedy podklad a delici linku; lze v ni sazet
% vice radku matematiky. Obe casti maji jen jeden spolecny ramecek.
\tcbset{solutionlesson/.style={
  commonlesson,
  enhanced jigsaw,
  breakable,
  colbacklower=black!3,
  segmentation style={solid,draw=black!25,line width=0.35pt},
  fontlower=\small\color{notegray},
  before lower={\noindent\textbf{Řešení:}\enspace},
  leftlower=3mm,rightlower=3mm
}}
\newtcolorbox{prikladReseniBox}[1][]{solutionlesson,title={Příklad},#1}
\newtcolorbox{zadaniReseniBox}[1][]{solutionlesson,title={Zadání},#1}

% Pomucky pro autora pripravy: sedive poznamky; \reseni pro starsi pripravy
\newcommand{\poznamka}[1]{\par\begingroup\small\color{notegray}\itshape
  \textnormal{\textbf{Poznámka pro učitele.}}\enspace #1\par\endgroup}
\newcommand{\reseni}[1]{\par\begingroup\small\color{notegray}
  \textbf{Očekávané řešení:}\enspace #1\par\endgroup}
\newcommand{\blok}[1]{\needspace{5\baselineskip}\subsection*{#1}}
\newcommand{\vyrok}[1]{\uv{#1}}

\usepackage{tikz}
\nastavPripravu{1. ročník}{Funkce}{72}{Opakování funkcí a přechod k rovnici}
\begin{document}
\blok{Od hodnoty funkce k otázce po argumentu}
Při výpočtu $f(2)$ známe argument a hledáme výstup. V rovnici $f(x)=5$ známe požadovaný výstup a hledáme všechny přípustné argumenty. Graficky jde o průsečíky grafu $f$ s vodorovnou přímkou $y=5$. Dnes tento vztah ukážeme na známých lineárních grafech; pravidla ekvivalentních úprav přijdou v hodině 73.
\begin{definiceBox}
\textbf{Řešení rovnice} $L(x)=P(x)$ na určeném oboru $D$ je každé $x\in D$, pro které mají oba výrazy po dosazení stejnou hodnotu. \textbf{Množina řešení} je množina všech takových argumentů.
\end{definiceBox}
\begin{prikladReseniBox}
Pro $f(x)=2x+1$ na $\mathbb R$ určete $f(0)$, nulovou hodnotu, argument s hodnotou $5$ a význam obou úloh v grafu.
\tcblower
$f(0)=1$. Nulová hodnota nastane, když $2x+1=0$, tedy při $x=-\frac12$; bod $(-\frac12;0)$ leží na ose $x$. Hodnota $5$ nastane při $2x+1=5$, tedy $x=2$ a průsečík s vodorovnou přímkou $y=5$ je $(2;5)$. Dosazením: $f(-\frac12)=0$ a $f(2)=5$.
\end{prikladReseniBox}
\blok{Rovnost dvou funkčních hodnot}
Z hodiny 68 známe porovnání dvou přímek. Rovnice $f(x)=g(x)$ se ptá, pro které stejné $x$ jsou body obou grafů ve stejné výšce. Výsledek musí patřit do společného definičního oboru.
\begin{prikladReseniBox}
Sestrojte $f(x)=x+2$ a $g(x)=-x+4$ na $\mathbb R$; z grafu odečtěte řešení $f(x)=g(x)$ a ověřte je výpočtem.
\tcblower
Pro $f$ použijeme body $(0;2),(2;4)$, pro $g$ body $(0;4),(2;2)$. Grafy se protínají v $(1;3)$. Početně $x+2=-x+4$, tedy $2x=2$ a $x=1$. Zkouška: $f(1)=3=g(1)$. Množina řešení je $K=\{1\}$; rovnice hledá první souřadnici průsečíku, nikoli celou dvojici.
\end{prikladReseniBox}
\begin{center}
\begin{tikzpicture}[x=.85cm,y=.7cm,>=stealth]
\draw[step=1,gray!25,very thin] (-1,-1) grid (4,5);
\draw[->] (-1.3,0)--(4.4,0) node[right] {$x$};\draw[->] (0,-1.2)--(0,5.4) node[above] {$y$};
\foreach \x in {1,2,3,4}{\draw (\x,.08)--(\x,-.08) node[below] {\small $\x$};}
\foreach \y in {1,2,3,4}{\draw (.08,\y)--(-.08,\y) node[left] {\small $\y$};}
\draw[very thick] (-1,1)--(3,5) node[above right] {\small $f$};
\draw[very thick,dashed] (-1,5)--(4,0);
\node[above right] at (3,1) {\small $g$};
\fill (1,3) circle (2.5pt) node[above right=5pt] {\small $(1;3)$};
\end{tikzpicture}
\end{center}
\poznamka{Nech žáky nejprve odečíst průsečík. Výpočet je kontrola přesné první souřadnice. Přímku vymezujeme jen oknem obrázku, nikoli skutečným omezením $D(f)=\mathbb R$.}
\newpage
\blok{Opakování v souvislostech}
\begin{zadaniReseniBox}
\begin{enumerate}[label=\alph*)]
\item Funkce $h(x)=-x+3$ je zadána na $\langle0;4\rangle$. Určete $D,H$, oba krajní body, průsečíky s osami a monotonii.
\item Pro $p(x)=|x-1|+2$ určete vrchol, obor hodnot a všechna $x$ s $p(x)=3$.
\item Dvě dílny mají náklady $A(n)=50n+300$ a $B(n)=70n+100$ Kč, kde $n\in\mathbb N_0$. Při jakém počtu kusů jsou náklady stejné? Která dílna je levnější při pěti kusech?
\end{enumerate}
\tcblower
\begin{enumerate}[label=\alph*)]
\item Krajní body $(0;3),(4;-1)$, tedy $D=\langle0;4\rangle$, $H=\langle-1;3\rangle$. Průsečíky jsou $(0;3)$ a $(3;0)$. Směrnice $-1$ znamená pokles.
\item Vrchol $(1;2)$, $H=\langle2;\infty)$. $|x-1|+2=3$ dává $|x-1|=1$, tedy $x=0$ nebo $x=2$; obě hodnoty kontrolou vyhovují.
\item $50n+300=70n+100$ znamená $200=20n$, tedy $n=10$. Obor $\mathbb N_0$ jej připouští. Při $n=5$ je $A(5)=550$ Kč a $B(5)=450$ Kč, takže levnější je B.
\end{enumerate}
\end{zadaniReseniBox}
\blok{Volitelné rozšíření a rezerva}
\begin{zadaniReseniBox}
Funkce $u(x)=2x-1$ je zadána jen pro $x\in\langle0;2\rangle$ a $v(x)=x+1$ na $\mathbb R$. Vyřešte $u(x)=v(x)$. Proč nestačí rovnost předpisů na celé $\mathbb R$?
\tcblower
$2x-1=x+1$ dává $x=2$, který patří do společného oboru $\langle0;2\rangle$; $u(2)=v(2)=3$. Kdyby algebraický kořen ležel mimo obor $u$, nebyl by řešením zadané rovnice funkcí.
\end{zadaniReseniBox}
\begin{zadaniReseniBox}
Načrtněte graf $y=|x|$ a odhadněte počet řešení $|x|=c$ pro $c=-1,0,2$.
\tcblower
Vodorovná přímka pod osou $x$ graf neprotne, tedy pro $c=-1$ žádné řešení. Při $c=0$ se dotkne vrcholu: $x=0$. Při $c=2$ protne obě větve: $x=-2,2$.
\end{zadaniReseniBox}
\blok{Domácí práce}
\begin{zadaniReseniBox}
\begin{enumerate}[label=\alph*)]
\item Určete z grafu a poté početně, kdy $2x+3=7$.
\item Najděte společnou hodnotu $f(x)=3x-2$ a $g(x)=x+4$ a příslušný argument.
\item Rozhodněte, zda rovnice $x+1=4$ má řešení pro funkci $f(x)=x+1$ na $\langle0;2\rangle$.
\end{enumerate}
\tcblower
\begin{enumerate}[label=\alph*)]
\item Přímky $y=2x+3$ a $y=7$ se protínají při $x=2$; kontrola $2\cdot2+3=7$.
\item $3x-2=x+4$ dává $2x=6$, $x=3$; společná hodnota je $7$ a průsečík $(3;7)$.
\item Předpisová rovnice dává $x=3$, ale $3\notin\langle0;2\rangle$. Pro zadanou funkci tedy řešení neexistuje.
\end{enumerate}
\end{zadaniReseniBox}
\poznamka{Jádro: motivace a vzor 9 min, grafická rovnost 14 min, samostatné opakování a kontrola 19 min, organizace 3 min. Rezerva nad rámec; domácí práce asi 15 min. Dnes používáme jen snadné rovnosti známých předpisů, úplná teorie ekvivalentních úprav patří hodině 73.}
\end{document}
